% 图 1.X：研究背景图——视觉基础模型驱动的开放视觉感知
% 风格参考：漏斗式三层结构（问题 → 挑战 → 方案）
% 简洁学术风格，不堆砌技术细节

\begin{tikzpicture}[
    x=1.0cm, y=1.0cm,
    every node/.style={font=\small},
    % 第一层：顶层标题
    toptitle/.style={rectangle, draw=none, fill=none,
                     font=\bfseries\small, align=center},
    % 第二层：三类开放性约束
    constraintbox/.style={rectangle, draw, very thick, rounded corners=4pt,
                          minimum width=4.5cm, minimum height=2.8cm, align=left,
                          inner sep=6pt, fill=white},
    domain/.style={constraintbox, draw=blue!60!black, fill=blue!6},
    category/.style={constraintbox, draw=green!50!black, fill=green!6},
    semantic/.style={constraintbox, draw=orange!60!black, fill=orange!6},
    % 缺陷标签
    defectlabel/.style={rectangle, draw, thick, rounded corners=3pt,
                        fill=yellow!20, draw=yellow!60!black,
                        font=\footnotesize, align=center},
    % 第三层：本文方案
    methodbox/.style={rectangle, draw, very thick, rounded corners=4pt,
                      minimum width=1.8cm, minimum height=1.5cm, align=center,
                      inner sep=4pt, font=\footnotesize},
    ch2/.style={methodbox, draw=blue!70!black, fill=blue!8},
    ch3/.style={methodbox, draw=teal!60!black, fill=teal!8},
    ch4/.style={methodbox, draw=green!60!black, fill=green!8},
    ch5/.style={methodbox, draw=orange!70!black, fill=orange!8},
    ch6/.style={methodbox, draw=red!60!black, fill=red!8},
    % 第四层：目标
    goalbox/.style={rectangle, draw, very thick, rounded corners=4pt,
                    minimum width=8.0cm, minimum height=1.2cm, align=center,
                    inner sep=6pt, fill=gray!8, draw=gray!50!black},
    % 连线样式
    topline/.style={->, >=Stealth, thick, draw=gray!60!black},
    arrowdown/.style={->, >=Stealth, very thick, draw=red!50!black},
    subarrow/.style={->, >=Stealth, thick, dashed, draw=gray!60!black},
    % 标注
    labellayer/.style={font=\small\bfseries, color=gray!40!black},
    notebox/.style={rectangle, draw, dashed, thick, rounded corners=3pt,
                    fill=yellow!15, draw=yellow!50!black,
                    font=\footnotesize\itshape, align=center}
]

% ========== 第一层：顶层标题 ==========
\node[toptitle] at (10.0, 11.5) {
    \textbf{研究背景：视觉基础模型驱动的开放视觉感知}
};

% ========== 第二层：三类开放性约束 ==========

% 左：域开放
\node[domain] at (3.5, 8.5) {
    \textbf{域开放}\\[3pt]
    \footnotesize 训练与测试数据分布偏移\\[4pt]
    \textbf{核心挑战：}\\[2pt]
    \begin{itemize}\setlength\itemsep{1pt}
        \item 伪标签噪声难以抑制
        \item 局部建模与长程依赖两难
    \end{itemize}
};

% 中：类别开放
\node[category] at (10.0, 8.5) {
    \textbf{类别开放}\\[3pt]
    \footnotesize 训练类别与测试类别不对称\\[4pt]
    \textbf{核心挑战：}\\[2pt]
    \begin{itemize}\setlength\itemsep{1pt}
        \item 单一语义锚泛化能力有限
        \item 多模态知识浅层融合
    \end{itemize}
};

% 右：语义开放
\node[semantic] at (16.5, 8.5) {
    \textbf{语义开放}\\[3pt]
    \footnotesize 低层感知向高层认知推理延伸\\[4pt]
    \textbf{核心挑战：}\\[2pt]
    \begin{itemize}\setlength\itemsep{1pt}
        \item 缺乏高阶情感推理能力
        \item 评测粒度不足
    \end{itemize}
};

% 三个约束下方的共同标注
\node[defectlabel] at (10.0, 6.5) {
    ◀ 现有方法在开放场景下存在泛化瓶颈 ▶
};

% ========== 第二层 → 第三层 连接箭头 ==========
\draw[arrowdown] (3.5, 6.8) -- (3.5, 5.8);
\draw[arrowdown] (10.0, 6.8) -- (10.0, 5.8);
\draw[arrowdown] (16.5, 6.8) -- (16.5, 5.8);

% ========== 第三层：本文方案 ==========

% 本文方案总标题
\node[labellayer] at (10.0, 5.5) {
    \textbf{本文方案：开放视觉感知渐进式解决}
};

% Ch2
\node[ch2] at (2.0, 4.0) {
    \textbf{第二章}\\[3pt]
    伪标签去噪\\[2pt]
    \itshape 域适应
};

% Ch3
\node[ch3] at (5.0, 4.0) {
    \textbf{第三章}\\[3pt]
    长程高效建模\\[2pt]
    \itshape 状态空间
};

% Ch4
\node[ch4] at (8.0, 4.0) {
    \textbf{第四章}\\[3pt]
    多模态协同\\[2pt]
    \itshape 知识蒸馏
};

% Ch5
\node[ch5] at (11.0, 4.0) {
    \textbf{第五章}\\[3pt]
    双先验融合\\[2pt]
    \itshape 开放词汇
};

% Ch6
\node[ch6] at (14.0, 4.0) {
    \textbf{第六章}\\[3pt]
    诊断评测\\[2pt]
    \itshape 语义推理
};

% 方案之间的递进箭头
\draw[subarrow] (3.3, 4.0) -- (4.0, 4.0);  % Ch2 → Ch3
\draw[subarrow] (6.0, 4.0) -- (6.8, 4.0);  % Ch3 → Ch4
\draw[subarrow] (9.0, 4.0) -- (9.8, 4.0);  % Ch4 → Ch5
\draw[subarrow] (12.0, 4.0) -- (13.0, 4.0); % Ch5 → Ch6

% 补充标注：约束与章节的对应关系
\draw[dashed, ->, >=Stealth, thick, draw=blue!60!black] (3.5, 7.1) .. controls (3.5, 6.2) and (2.5, 5.0) .. (2.5, 4.6);
\draw[dashed, ->, >=Stealth, thick, draw=green!50!black] (10.0, 7.1) .. controls (10.0, 6.2) and (9.5, 5.0) .. (9.5, 4.6);
\draw[dashed, ->, >=Stealth, thick, draw=orange!60!black] (16.5, 7.1) .. controls (16.5, 6.2) and (14.5, 5.0) .. (14.5, 4.6);

% ========== 第四层：研究目标 ==========
\draw[arrowdown] (10.0, 3.0) -- (10.0, 2.2);

\node[goalbox] at (10.0, 1.6) {
    \textbf{研究目标}\\[4pt]
    \footnotesize 构建基础模型驱动的开放视觉感知方法体系\\
    \itshape 精度提升 · 效率优化 · 可解释性增强
};

% ========== 辅助标注：开放性维度递进 ==========
\node[notebox] at (18.0, 8.5) {
    \textbf{开放性维度}\\[3pt]
    域开放 → 类别开放\\
    → 语义开放
};

\end{tikzpicture}
